In multi-environment causal representation learning, environment
information helps separate stable structure from environment-specific
variation. Whether a downstream analysis holds up depends on whether the
relevant environment structure can be identified or reasonably
approximated from unlabeled data. This paper studies that preceding
problem, not causal representation learning itself. We ask what latent
structure is identifiable from pooled observations and when
recording-level clustering can recover groups under explicit
assumptions, with a focus on a covariance-based environment structure
once the labels are removed.

Most multi-environment theory takes the partition as given. The index
set of each environment is known and unique, so identification and
estimation proceed with the assignment held fixed. This is convenient in
designed experiments, where interventions are logged and the source of
every observation is recorded. It becomes less secure once the labels
disappear, because a pooled marginal, written \(P_X=\sum_e\pi_e G_e\),
carries no mark of the component behind any given draw. Missing labels
do not, on their own, settle the question. Three issues need to be kept
apart. The labels themselves are unobserved; whether the component
parameters are uniquely determined by the distribution is a separate
matter; and whether the true label of an individual observation can be
recovered without error is a third. Under a restricted model the mixture
parameters may still be identifiable, and even when they are,
overlapping components can leave individual labels impossible to judge
with certainty. Parameter identification turns on the model's
constraints, while label recovery also depends on the observation unit,
the repeated measurements, and whether the clustering objective lines up
with the latent grouping.

Two questions organize the analysis. First, once only a pooled
distribution remains, is the covariance-based environment structure
identifiable, and where does the boundary between non-identifiability
and point identification lie? Second, under what explicit conditions can
recording-level clustering recover latent groups, and how do
cluster-number selection and the sampling design affect empirical
recovery? The first asks what the population distribution can determine
and is taken up in Section 4. The second asks how a recording-level
procedure behaves under added structural assumptions; its theory is
developed in Section 5 and its empirical evidence in Section 7. The two
are related but distinct.

One central claim organizes the answer. Under unlabeled pooling,
covariance-based environment discovery has a sharp identifiability
boundary and a conditional empirical recovery boundary:
source-associated technical structure is recoverable at a fixed number
of environments, while blind selection of that number and any causal
reading of the recovered groups remain unreliable.

Three lines of prior work frame the gap, and the full four-category
review is left to Section 2. Multi-environment representation learning
and environment inference either assume the environment is known or
build a partition that serves a prediction target, so they do not say
what is identifiable, or under what conditions grouping can be
recovered, once the labels are absent. Mixture identifiability recovers
component parameters from the population distribution or its moments,
yet it treats parameter identification and recording-level label
recovery as one object, and the two are not the same. The K-means and
covariance-geometry literature establishes the statistical properties of
the clustering objective and distance representations of covariance
features, but it does not settle when that objective aligns with the
target environment grouping or where the guarantees of a deployed
pipeline end. The gap is not that a particular construction is missing.
These results, framed for known environments, parameter recovery, or a
generic objective, simply do not answer the two questions above.

The motivation follows from the distance between two inferences that are
easy to conflate. Identifiable mixture parameters do not make recording
labels recoverable, and a stable clustering that tracks the source
labels does not make the recovered groups causal environments. Without
keeping these apart, a clustering result can be assigned more meaning
than the evidence supports in a downstream analysis. The concrete route
into the problem is the repeated recording. A single pooled observation
offers little beyond the marginal, but a recording contains several
windows, and their within-recording moments can be aggregated into a
recording-level covariance feature. This turns group recovery into a
well-defined statistical problem, and it also introduces the conditions
that govern it: recording length, the accuracy of the feature estimate,
and the choice of analysis unit.

At the population level, a concrete \(d=2\) construction gives two
partitions that induce the same pooled distribution, disagree on half
the support, and are not related by a label permutation. The witness
shows that identification cannot be granted without constraints, the
expected state for arbitrary measurable partitions. The contrast is a
constrained result. In a two-dimensional, equally weighted Gaussian
mixture with symmetric covariance deviations around the identity,
fourth-order cumulants identify the unordered component covariance pair.
Relaxing the setting leaves a family of structures consistent with the
pooled data, a region we discuss without claiming a sharp set or a
coverage result.

At the recording level, the DELVE procedure builds a feature model and a
finite-sample risk-transfer argument. Closed level sets and compactness
give a positive risk gap, and the empirical excess risk is then
converted into a control on the grouping error, with population
alignment stated as an explicit compatibility premise. The empirical
study applies averaged log-covariance features across three EEG
datasets. On a label-assisted balanced design of 90 recordings,
known-\(K\) clustering reached a recording-level adjusted Rand index
(ARI) of \(0.582\pm0.068\) across six prespecified seeds. The automatic
selector chose two clusters in three seeds and three in the other three
and reached \(0.419\pm0.140\). The unbalanced and subject-level plug-in
analyses were weaker, at \(0.371\) and \(0.241\), and equalizing the
window count gave a similar ARI of \(0.577\), so window count alone does
not account for the grouping. The reading is deliberately narrow. In
this balanced design, known-\(K\) source grouping shows some recovery,
and the real-data label mixes device, task, population, duration, and
preprocessing; it is not a biological or causal environment.

The standard components carry no novelty. Shrinkage covariance
estimation, the log-Euclidean feature, nearest-centroid clustering, and
the resampling diagnostics are established tools, each cited for the
exact construction used. The three contributions below are the
components of the one claim above, not independent results of equal
strength.

\begin{enumerate}
\def\labelenumi{\arabic{enumi}.}
\item
  \textbf{Boundary frame.} Through an explicit witness and a constrained
  fourth-order identification, we separate the non-uniqueness of
  unstructured partitions from the identifiability of structured
  component parameters, with the relaxed case treated as a discussion
  and not a sharp set.
\item
  \textbf{Theoretical pillar.} We define the recording-level estimand
  and build the risk-transfer route, from a positive risk gap under
  closed level sets and compactness to a control that converts empirical
  excess risk into grouping error, with the applicable object and the
  population-alignment premise stated explicitly.
\item
  \textbf{Empirical pillar.} Across three independent EEG datasets, we
  treat known-\(K\) grouping performance, the reliability of automatic
  \(K\) selection, and the causal meaning of the source groups as three
  questions evaluated separately, supported by a pre-specified,
  unit-aware protocol and the fixed-window, unbalanced, and
  within-source controls.
\end{enumerate}

The stance stays upstream of causal representation learning, which
remains a potential downstream use and is not a task completed here.
Section 2 reviews the related work and Section 3 fixes the notation.
Section 4 sets the population boundary and Section 5 studies recovery
under a different observation mechanism, so the two are complementary
and not a theorem with its implementation. Sections 6 and 7 report the
synthetic and EEG studies under their own designs, and Section 8 states
the limits.
